%========================================================
% APPENDIX A.2 - Derivation of Fidelity for a Pure Reference State
%========================================================

\section{Derivation of Fidelity for a Pure Reference State}
\label{app:fidelity_pure_reference_state}

This appendix derives the simplified expression for quantum fidelity when one of the two states is pure. The result is used throughout Chapter~4 in the evaluation of fidelity with respect to Bell states.

\medskip

Beyond its algebraic simplification, this derivation also clarifies the operational interpretation of fidelity as a projection probability onto a target quantum state.

%--------------------------------------------------------
% General definition of fidelity
%--------------------------------------------------------

\subsection*{General Definition of Fidelity}

Let \( \rho \) and \( \sigma \) be density operators acting on a finite-dimensional Hilbert space.

The general definition of quantum fidelity is

\begin{equation}
F(\rho,\sigma)
=
\left(
\operatorname{Tr}
\sqrt{
\sqrt{\rho}\,
\sigma\,
\sqrt{\rho}
}
\right)^2.
\label{eq:general_fidelity_definition}
\end{equation}

In the present derivation, we consider the special case in which the reference state is pure:
\begin{equation}
\sigma
=
|\psi\rangle\langle\psi|,
\label{eq:pure_reference_density_operator}
\end{equation}
with
\begin{equation}
\langle \psi | \psi \rangle = 1.
\label{eq:pure_state_normalization}
\end{equation}

By construction, the operator \( \sigma \) is:

\begin{itemize}
    \item Hermitian,
    \[
    \sigma^\dagger = \sigma,
    \]

    \item positive semidefinite,
    \[
    \langle \chi | \sigma | \chi \rangle \ge 0,
    \qquad
    \forall\,|\chi\rangle,
    \]

    \item normalized,
    \[
    \operatorname{Tr}(\sigma)=1.
    \]
\end{itemize}

Moreover, \( \sigma \) is a rank-one projector onto the one-dimensional subspace spanned by \( |\psi\rangle \).

%--------------------------------------------------------
% Reduction for a pure reference state
%--------------------------------------------------------

\subsection*{Reduction for a Pure Reference State}

Substituting Eq.~\eqref{eq:pure_reference_density_operator} into Eq.~\eqref{eq:general_fidelity_definition} gives

\begin{equation}
\sqrt{\rho}\,
\sigma\,
\sqrt{\rho}
=
\sqrt{\rho}\,
|\psi\rangle\langle\psi|\,
\sqrt{\rho}.
\label{eq:fidelity_intermediate_operator}
\end{equation}

Define the vector

\begin{equation}
|\phi\rangle
=
\sqrt{\rho}\,
|\psi\rangle.
\label{eq:phi_definition}
\end{equation}

Then Eq.~\eqref{eq:fidelity_intermediate_operator} becomes

\begin{equation}
\sqrt{\rho}\,
\sigma\,
\sqrt{\rho}
=
|\phi\rangle\langle\phi|.
\label{eq:rank_one_operator_phi}
\end{equation}

The operator \( |\phi\rangle\langle\phi| \) is again rank one and positive semidefinite.

%--------------------------------------------------------
% Spectral structure of the rank-one operator
%--------------------------------------------------------

\subsection*{Spectral Structure of the Rank-One Operator}

The only nonzero eigenvalue of the operator in Eq.~\eqref{eq:rank_one_operator_phi} is

\begin{equation}
\lambda
=
\langle \phi | \phi \rangle.
\label{eq:rank_one_eigenvalue}
\end{equation}

Introducing the normalized vector

\begin{equation}
|\hat{\phi}\rangle
=
\frac{|\phi\rangle}{\|\phi\|},
\label{eq:normalized_phi}
\end{equation}

the spectral decomposition becomes

\begin{equation}
|\phi\rangle\langle\phi|
=
\lambda\,
|\hat{\phi}\rangle\langle\hat{\phi}|.
\label{eq:spectral_decomposition_rank_one}
\end{equation}

The square root of the operator is therefore

\begin{equation}
\sqrt{
|\phi\rangle\langle\phi|
}
=
\sqrt{\lambda}\,
|\hat{\phi}\rangle\langle\hat{\phi}|.
\label{eq:square_root_rank_one}
\end{equation}

Taking the trace gives

\begin{equation}
\operatorname{Tr}
\sqrt{
|\phi\rangle\langle\phi|
}
=
\sqrt{\lambda}.
\label{eq:trace_square_root_rank_one}
\end{equation}

Substituting into Eq.~\eqref{eq:general_fidelity_definition}, one obtains

\begin{equation}
F(\rho,\sigma)
=
\lambda
=
\langle \phi | \phi \rangle.
\label{eq:fidelity_lambda_relation}
\end{equation}

Using Eq.~\eqref{eq:phi_definition},

\begin{equation}
\lambda
=
\langle \psi |
\rho
|\psi\rangle.
\label{eq:lambda_final_expression}
\end{equation}

%--------------------------------------------------------
% Final expression and interpretation
%--------------------------------------------------------

\subsection*{Final Expression and Interpretation}

Combining the previous results, the fidelity simplifies to

\begin{equation}
F(\rho,|\psi\rangle\langle\psi|)
=
\langle \psi |
\rho
|\psi\rangle.
\label{eq:fidelity_pure_reference_final}
\end{equation}

This is the expression used throughout Chapter~4 for fidelity with respect to Bell states.

\medskip

The result admits a direct operational interpretation. Since
\[
|\psi\rangle\langle\psi|
\]
is the projector onto the state \( |\psi\rangle \), the fidelity corresponds to the probability of obtaining the outcome \( |\psi\rangle \) in a projective measurement performed on the state \( \rho \).

\medskip

Fidelity therefore quantifies the overlap between the physical state produced by the system and the chosen target reference state.